% !TeX program = xelatex
% !TeX encoding = UTF-8
\documentclass[12pt,a4paper]{article}
\usepackage{fontspec}
\usepackage[top=2.54cm,bottom=2.54cm,left=3.18cm,right=3.18cm,headheight=0pt,headsep=0pt,footskip=1.5cm]{geometry}
\setmainfont{NotoSerifSC-Regular.ttf}[Path=build-fonts/,BoldFont=NotoSerifSC-Bold.ttf]
\newfontfamily\headingfont{NotoSansSC-Regular.ttf}[Path=build-fonts/]
\newcommand{\songti}{\normalfont}
\newcommand{\heiti}{\headingfont}
\XeTeXlinebreaklocale "zh"
\XeTeXlinebreakskip=0pt plus 1pt
\usepackage{amsmath,amssymb,mathtools,bm}
\renewcommand{\theequation}{\arabic{equation}}
\usepackage{graphicx}
\usepackage{longtable,booktabs,array,tabularx,calc}
\usepackage{etoolbox}
\usepackage{float}
\usepackage{needspace}
\usepackage{caption}
\usepackage{setspace}
\usepackage{fancyhdr}
\usepackage{titlesec}
\usepackage{tocloft}
\usepackage{indentfirst}
\usepackage{enumitem}
\usepackage{fvextra}
\usepackage{xcolor}
\usepackage{xurl}
\usepackage[hidelinks]{hyperref}
\usepackage{bookmark}
\usepackage{eso-pic}
\usepackage{tikz}
\usetikzlibrary{positioning,arrows.meta,shapes.geometric,fit,backgrounds,calc}
\newcommand{\R}{\mathbb{R}}
\newcommand{\Bset}{\{0,1\}}
\DeclareMathOperator*{\argmin}{arg\,min}
\definecolor{paperGreen}{RGB}{132,168,60}
\definecolor{paperGreenDark}{RGB}{79,96,48}
\definecolor{paperGreenLight}{RGB}{212,224,192}
\definecolor{paperCream}{RGB}{245,243,232}
\tikzset{
  startstop/.style={rectangle,rounded corners=8pt,minimum width=3.2cm,minimum height=0.9cm,align=center,draw=paperGreenDark,fill=paperGreen,text=white,thick},
  process/.style={rectangle,minimum width=3.6cm,minimum height=0.9cm,align=center,draw=paperGreenDark,fill=paperCream,thick},
  decision/.style={diamond,aspect=2.1,minimum width=3.0cm,align=center,draw=paperGreenDark,fill=paperGreenLight,thick,inner sep=1pt},
  data/.style={trapezium,trapezium left angle=70,trapezium right angle=110,minimum width=3.2cm,minimum height=0.9cm,align=center,draw=paperGreenDark,fill=paperGreenLight!55,thick},
  verify/.style={rectangle,minimum width=3.4cm,minimum height=0.9cm,align=center,draw=paperGreenDark,fill=paperGreenLight,thick},
  layer/.style={rectangle,draw=paperGreenDark,dashed,inner sep=6pt},
  arrow/.style={-{Stealth[length=2.4mm]},thick,draw=paperGreenDark},
  lab/.style={font=\small,midway}
}
\newcommand{\sanhao}{\fontsize{16pt}{16pt}\selectfont}
\newcommand{\xiaosan}{\fontsize{15pt}{18.75pt}\selectfont}
\newcommand{\sihao}{\fontsize{14pt}{14pt}\selectfont}
\newcommand{\xiaosi}{\fontsize{12pt}{12pt}\selectfont}
\newcommand{\wuhao}{\fontsize{10.5pt}{10.5pt}\selectfont}
% 正文参考文献编号采用右上角标注
\newcommand{\supcite}[1]{\textsuperscript{[#1]}}

\AtBeginDocument{%
  \xiaosi\songti\setstretch{1.25}%
  % 统一所有行间公式与上下正文之间的距离
  \setlength{\abovedisplayskip}{7.5pt}%
  \setlength{\belowdisplayskip}{7.5pt}%
  \setlength{\abovedisplayshortskip}{7.5pt}%
  \setlength{\belowdisplayshortskip}{7.5pt}%
}
\setstretch{1.25}
\setlength{\parindent}{2em}
\setlength{\parskip}{7.5pt}
\pagestyle{fancy}
\fancyhf{}
\fancyfoot[C]{\xiaosi\thepage}
\renewcommand{\headrulewidth}{0pt}
\renewcommand{\footrulewidth}{0pt}
\fancypagestyle{plain}{\fancyhf{}\fancyfoot[C]{\xiaosi\thepage}\renewcommand{\headrulewidth}{0pt}\renewcommand{\footrulewidth}{0pt}}
\newcommand{\ChineseNumber}[1]{%
  \ifcase#1 零\or 一\or 二\or 三\or 四\or 五\or 六\or 七\or 八\or 九\or 十%
  \or 十一\or 十二\or 十三\or 十四\or 十五\or 十六\or 十七\or 十八\or 十九\or 二十\else #1\fi}
\renewcommand{\thesection}{\ChineseNumber{\value{section}}、}
\renewcommand{\thesubsection}{（\ChineseNumber{\value{subsection}}）}
\renewcommand{\thesubsubsection}{（\arabic{subsubsection}）}
\renewcommand{\theparagraph}{（\arabic{paragraph}）}
\renewcommand{\thesubparagraph}{\arabic{subparagraph}）}
\titleformat{\section}{\centering\heiti\xiaosan}{\thesection}{0.5em}{}
\titlespacing*{\section}{0pt}{12pt}{6pt}
\titleformat{\subsection}{\raggedright\heiti\sihao}{\thesubsection}{0.5em}{}
\titlespacing*{\subsection}{0pt}{6pt}{0pt}
\titleformat{\subsubsection}{\raggedright\songti\xiaosi\bfseries}{\thesubsubsection}{0.5em}{}
\titlespacing*{\subsubsection}{0pt}{6pt}{0pt}
\titleformat{\paragraph}[runin]{\raggedright\songti\xiaosi}{\theparagraph}{0.5em}{}
\titleformat{\subparagraph}[runin]{\raggedright\songti\xiaosi}{\thesubparagraph}{0.5em}{}
\setcounter{secnumdepth}{5}
\setcounter{tocdepth}{3}
% 目录按标题层级逐级缩进：一级顶格，二级和三级依次右移
\cftsetindents{section}{0em}{4em}
\cftsetindents{subsection}{0em}{4em}
\cftsetindents{subsubsection}{0em}{4em}
\makeatletter
\let\paper@l@subsection\l@subsection
\let\paper@l@subsubsection\l@subsubsection
\renewcommand*{\l@subsection}[2]{\paper@l@subsection{\hspace*{2em}#1}{#2}}
\renewcommand*{\l@subsubsection}[2]{\paper@l@subsubsection{\hspace*{4em}#1}{#2}}
\makeatother
\renewcommand{\figurename}{图}
\renewcommand{\tablename}{表}
\DeclareCaptionFont{captioncn}{\songti\xiaosi\setstretch{1.0}}
\DeclareCaptionLabelSeparator{twospaces}{\hspace{2em}}
\captionsetup{font=captioncn,labelsep=twospaces,justification=centering,singlelinecheck=true}
\captionsetup[table]{position=top}
\captionsetup[figure]{position=bottom}
\renewcommand{\thefootnote}{\textcircled{\scriptsize\arabic{footnote}}}
\graphicspath{{media/}}
\makeatletter
\def\maxwidth{\ifdim\Gin@nat@width>\linewidth\linewidth\else\Gin@nat@width\fi}
\def\maxheight{\ifdim\Gin@nat@height>0.82\textheight 0.82\textheight\else\Gin@nat@height\fi}
\def\fps@figure{htbp}
\patchcmd\longtable{\par}{\if@noskipsec\mbox{}\fi\par}{}{}
\makeatother
\AtBeginEnvironment{tabular}{\songti\xiaosi\setstretch{1.0}}
\AtBeginEnvironment{longtable}{\songti\xiaosi\setstretch{1.0}}
\setkeys{Gin}{width=\maxwidth,height=\maxheight,keepaspectratio}
\setlength{\emergencystretch}{3em}
\providecommand{\tightlist}{\setlength{\itemsep}{0pt}\setlength{\parskip}{0pt}}
\urlstyle{same}
\hypersetup{pdfauthor={},pdftitle={暖城向新·热力不息}}

\newcommand{\EnableBodyBackground}{%
  \AddToShipoutPictureBG{%
    \AtPageLowerLeft{%
      \includegraphics[width=\paperwidth,height=\paperheight]{body-background.jpg}%
    }%
  }%
}

\newenvironment{refsectionflush}
  {\par\sloppy\setlength{\parindent}{0pt}\setlength{\parskip}{0pt}}
  {\par\fussy}
\newcommand{\refentry}[2]{%
  \par\noindent\hangindent=2em\hangafter=1
  [#1]\enspace #2\par}

\begin{document}

\AddToShipoutPictureBG*{%
  \AtPageLowerLeft{\includegraphics[width=\paperwidth,height=\paperheight]{cover.jpg}}%
}
\thispagestyle{empty}
\null
\clearpage
\EnableBodyBackground
\pagestyle{empty}
\begin{center}
{\songti\sanhao\bfseries\setstretch{1.0} 暖城向新·热力不息\par}
\vspace{18pt}
{\heiti\sihao 摘\quad 要\par}
\end{center}
\vspace{3pt}
\par\indent 城镇供热转型需要同时协调热源配置、管网输配、建筑用热、储热和配电能力。本文建立可按场景配置的源网荷储协同规划方法，保留既有集中供热重构、分散清洁供热和夏热冬冷地区新建系统的接口，以第三类场景原型组织完整冷热服务年及供暖季的软件验证。

建筑负荷采用精确离散的1R1C热状态方程，由节点交付的供热、供冷与室温共同决定；容量、温位、冷热互斥、储热连续和可靠性约束共同限定投资与调度。成本与系统总碳采用统一台账，多目标方法包含带容差记录的两阶段词典序和碳上限约束；代表期规划后冻结投资，以完整时序检验方案。

结果准入同时要求输入来源可核查、模型实际求解和独立数值审计。公开气象产品、观测数据、来源明确的工程参数与情景假设分别记录；无出处的建筑、价格或空间参数仅用于程序验证。气象来源复核、N2实测比较及原作者多目标例题复现形成三类证据：气象字段与官方响应一致，N2日块训练改善未转化为留出开环改善，五个数学采样点与解析界一致。上述证据分别限定适用范围，工程参数仍需真实资料闭合。

\vspace{-4.5pt}
\noindent{\heiti\xiaosi 关键词：}{\songti\xiaosi 城镇供热；源荷匹配；规划寻优；低碳转型；供热可靠性}
\clearpage
\thispagestyle{empty}
\begin{center}{\heiti\sihao 目\quad 录}\end{center}
\begingroup
\setlength{\parskip}{0pt}
\setlength{\cftbeforesecskip}{1pt}
\makeatletter
\@starttoc{toc}
\makeatother
\endgroup
\clearpage
\pagenumbering{arabic}
\setcounter{page}{1}
\pagestyle{fancy}
\section{赛题、场景与研究范围}

供热转型的规划对象包括热源、网络、建筑末端、储热和电网接入。仅按设计负荷配置单项设备，难以反映温位适配、逐时电力峰值、储热持续供能和建筑热惯性之间的约束。赛事要求模型覆盖源、网、荷、储，采用逐时尺度完成完整供暖季仿真，并分析经济性、碳排放与运行调度\supcite{1}。

三类场景共享守恒与资产接口。第一类侧重既有集中供热系统的资产保留、热源退役、余热接入及网络降温；第二类侧重分散设施、用户分配、燃料与微网边界；第三类侧重系统形态、冷热服务、建筑品质和配电容量。本稿以第三类场景原型组织软件验证，第一、二类保留适配范围。

模型分为城市配置、区域能流和建筑状态三个尺度。城市层确定设备与容量，区域层连接热源、管段、储热和配电，建筑层将交付冷热量映射为热状态。空间节点与建筑原型的真实性由输入资料决定，聚合拓扑及合成原型不能称为真实GIS或实测城市网络。

\begin{figure}[H]
\centering
\begin{tikzpicture}[font=\small]
\node[process,minimum width=3.6cm,text width=3.1cm] (city) at (0,0) {城市层\\容量与技术配置};
\node[process,minimum width=3.6cm,text width=3.1cm] (region) at (4.5,0) {区域层\\冷热、电力与储热};
\node[process,minimum width=3.6cm,text width=3.1cm] (building) at (9,0) {建筑层\\RC状态与服务边界};
\draw[{Stealth[length=2mm]}-{Stealth[length=2mm]},thick,draw=paperGreenDark] (city) -- (region);
\draw[{Stealth[length=2mm]}-{Stealth[length=2mm]},thick,draw=paperGreenDark] (region) -- (building);
\end{tikzpicture}
\caption{投资、能流与建筑状态的多尺度接口示意}
\label{fig:compact-multiscale}
\end{figure}

工程验证需通过数据来源、物理和优化模型、实际求解、独立审计四层证据。程序验证检查方程与实现，未经来源核实的原型结果限于“程序验证/情景示例”。本稿的工程经济与碳评价保留方法和核实因子，缺乏真实建筑、设备报价及公平基准时不计算工程改善率。历史三场景记录单独归档。

\clearpage
\section{数据来源、假设与准入}

\subsection{字段身份及追溯}

每个输入字段记录名称、数值或时序、单位、来源主体、原始链接或文件、获取时间、适用区域、换算过程及SHA256。数据性质至少区分公开产品、原始观测、来源明确的工程参数和情景假设。可访问一个网站不等于核实了具体数值；来源记录须能定位实际字段、年份与处理步骤。

\begin{table}[H]
\centering
\caption{输入字段与证据准入}
\begin{tabularx}{\textwidth}{@{}p{0.16\textwidth}p{0.27\textwidth}X@{}}
\toprule
类别 & 用途 & 工程结果准入要求\\
\midrule
气象 & 温度、辐照及湿冷性能输入 & 记录公开产品、坐标、时区、原始时间范围与下载文件；格点产品不表述为地面站实测\\
建筑 & 面积、热损、热容、得热、占用与舒适边界 & 定位设计资料或真实观测；无具体来源的原型参数只用于程序验证，不宣称已校准\\
空间资产 & 节点、候选管段、资源与接入容量 & 核对实际位置、距离、既有资产及候选形成依据；合成坐标和拓扑单列\\
设备价格 & 性能、投资、运维、电价与燃料 & 说明参数条件、报价或公开依据；未经核验的价格只作情景假设\\
政策与碳 & 服务范围、排放因子、碳价和可靠性 & 区分适用标准、统计因子和假设碳价；保留原口径，避免换算重复计量\\
\bottomrule
\end{tabularx}
\end{table}

建筑参数的不确定区间不能替代参数来源。H/C、太阳及内部得热、占用日程、电价和空间字段需分别审查；即使某组数据已通过求解、RC残差和费用台账核验，来源未明确时仍不得晋级为真实工程案例。历史材料适用相同门槛，不因曾写入报告而豁免。

气象采用NASA POWER逐时产品\supcite{9}。本次以API v2.10.2复核三城市2020--2024年共15份资料，131544小时的温度、相对湿度、10~m风速和短波辐射共526176个值与官方响应一致，原CSV未改。辐射实际为Wh/m$^2$，不能沿用旧字段名中的kWh解释；产品时钟为LST，转换须保留其定义。2023年旧记录没有原查询链接，本次重建查询单独留痕，不倒填旧获取日期或API版本。该复核仅认证四个产品字段，不认证生成负荷、建筑、价格或地面站测量精度。

\subsection{时间轴与单位}

原始时间戳保留UTC或供应商时区，转为项目本地时区后按左闭右开区间筛选。逐时输入检查唯一性、连续性、缺失和重复；完整服务年覆盖冷热服务，供暖季按服务合同单独筛选，不能只按论文中写出的天数推算结果时长。代表期的行数、小时权重和实际连续时长分别报告。

运行功率采用kW，能量采用kWh，时间步长采用h，排放采用kg；汇总使用MW、MWh和t时明确换算。热损系数为kW/K，热容为kWh/K。摄氏与开尔文温差数值相同，火用或温度比值涉及的绝对温度需另作换算。预处理不改变原始观测身份，插补及代理字段须标记并说明用途。

状态采样时刻与观测平均窗口分别登记。区间平均温度不能直接当作时段末温度，热状态和观测方程应按各自采样含义建立；时间序列验证还须检查残差与独立留出表现\supcite{3}。

\subsection{外部观测与留出验证}

外部RC观测记录许可、对象、测点、采样时刻和参数来源，参数辨识与留出评估使用分离的数据。误差由实际观测与模型输出对齐后计算；合成轨迹和模型自洽残差不能充当外部观测。本稿纳入N2首轮有限适用性诊断，来源审核完成的真实工程效果仍需独立资料。

\clearpage
\section{源网荷储与建筑热状态模型}

\subsection{变量、资产与多尺度接口}

令$b,n,k,e,s,t$分别表示建筑、区域节点、设备、管段、储热装置和时段。投资变量确定设备新增及退役容量、候选网络、储热和接入容量；运行变量包括交付冷热量、设备出力、节点能流、电力及热状态。设备可用容量写为
\begin{equation}
K_k=K_k^{\mathrm{exist}}+K_k^{\mathrm{new}}-K_k^{\mathrm{ret}},\qquad
0\leq Q_{k,t}\leq a_{k,t}K_k.
\label{eq:capacity}
\end{equation}
启停、最小出力、爬坡及政策退役条件按实际启用设备加载。城市投资变量通过容量上界进入区域运行，建筑交付量按节点与温位聚合后进入能量守恒。该接口保留各层身份，不能将规划投资间隙与固定投资后的调度间隙混为一项。

\subsection{建筑负荷与精确1R1C离散}

热状态结构参考建筑RC灰箱建模及模型适用性检验原则\supcite{2}。本文采用单状态简化，具体H/C及得热仍须依据对应建筑资料辨识，方法文献不提供武汉参数。

建筑以温度$\theta_b$为状态，热损系数$H_b>0$、热容$C_b>0$及得热$G_b$描述围护、换气与内部和太阳作用。在每小时内输入保持的条件下，热状态方程为
\begin{equation}
C_b\frac{\mathrm d\theta_b}{\mathrm dt}
=H_b(T^{\mathrm{out}}-\theta_b)+Q_b^{\mathrm h}-Q_b^{\mathrm c}+G_b.
\label{eq:rc-continuous}
\end{equation}
其中$Q_b^{\mathrm h},Q_b^{\mathrm c}$分别为交付供热及供冷功率，$G_b$为内部与太阳得热之和，采用kW；$C_b$采用kWh/K、$H_b$采用kW/K，时间采用h。取$\Delta t=1~\mathrm h$，定义
\begin{equation}
a_b=\exp\!\left(-\frac{H_b\Delta t}{C_b}\right),\qquad
\beta_b=\frac{1-a_b}{H_b}.
\label{eq:rc-coefficients}
\end{equation}
对式\eqref{eq:rc-continuous}积分得到
\begin{equation}
\theta_{b,t}=a_b\theta_{b,t-1}+(1-a_b)T_t^{\mathrm{out}}
+\beta_b\left(Q_{b,t}^{\mathrm h}-Q_{b,t}^{\mathrm c}+G_{b,t}\right).
\label{eq:alg-building}
\end{equation}
程序使用稳定的指数差计算系数，拒绝非正热损、热容及时间步长。供热与供冷变量直接参与上述状态和节点能流，不再叠加一套虚拟建筑储热需求，避免重复计量。

状态满足随占用及服务合同变化的边界，冷热互斥通过二元变量$m_{b,t}$限定：
\begin{align}
\underline\theta_{b,t}\leq\theta_{b,t}&\leq\overline\theta_{b,t},\\
0\leq Q_{b,t}^{\mathrm h}&\leq M_{b,t}^{\mathrm h}m_{b,t},\\
0\leq Q_{b,t}^{\mathrm c}&\leq M_{b,t}^{\mathrm c}(1-m_{b,t}).
\label{eq:rc-exclusive}
\end{align}
$M$为经输入检查计算的有限上界；不具有两种服务能力的时段按对应边界处理。完整时序采用声明的初始及终端状态，代表期按其组边界使用固定或循环条件。逐日固定同一初末温度会改变冷热需求，须单独核验其影响，不能以权重覆盖完整年度代替状态连续性。

\subsection{区域冷热平衡与温位}

节点$n$、温位$g$的热量守恒为
\begin{multline}
\sum_kQ_{k,n,g,t}+\sum_{e\in\delta^-(n)}\eta_eF_{e,g,t}
+\sum_sQ_{s,n,g,t}^{\mathrm{dis}}\\
=\sum_{b\in\mathcal B_{n,g}}Q_{b,t}^{\mathrm h}
+\sum_sQ_{s,n,g,t}^{\mathrm{ch}}+\sum_{e\in\delta^+(n)}F_{e,g,t}.
\label{eq:node-heat-balance}
\end{multline}
供冷节点按交付冷量建立对应平衡；已启用生活热水或其他服务时，将其独立需求加入相应账本。禁止用全系统总量平衡掩盖局部输配不足，也不将未满足服务人为删除以降低碳排放。

直接热源与用户温位的兼容要求为
\begin{equation}
T_k^{\mathrm{src,min}}\geq T_g^{\mathrm{sup}}+\Delta T^{\mathrm{hx}}.
\label{eq:alg-grade}
\end{equation}
低温资源经热泵提升后按其供热侧温位接入；式\eqref{eq:alg-grade}不用于将低温热泵资源端误当作直接供热端。管段容量与候选状态限制热流，连续温度、水力和施工可行性另作专业复核。本稿不把候选聚合管段称为已认证的实际管网。

\subsection{设备、储热与电力}

锅炉、热泵及制冷设备的转换采用输入确定的效率系数：
\begin{equation}
Q_{k,t}^{\mathrm{boiler}}=\eta_kF_{k,t}^{\mathrm{fuel}},\qquad
Q_{k,t}^{\mathrm{HP}}=COP_{k,t}P_{k,t}^{\mathrm{HP}},\qquad
Q_{k,t}^{\mathrm{cool}}=EER_{k,t}P_{k,t}^{\mathrm{cool}}.
\end{equation}
系数的温度、湿冷和部分负荷适用条件需记录来源；外生性能修正不等于建立了室内热湿或声学动态模型。冷热共用设备同时受投资容量与运行模式限制。

储热状态为
\begin{equation}
E_{s,t}=(1-\mu_s)E_{s,t-1}+\eta_s^{\mathrm{ch}}Q_{s,t}^{\mathrm{ch}}\Delta t
-\frac{Q_{s,t}^{\mathrm{dis}}\Delta t}{\eta_s^{\mathrm{dis}}},
\label{eq:storage-state}
\end{equation}
并满足能量容量、充放功率、互斥和初末边界。节点电力平衡将基础负荷、热泵、制冷、循环泵及已启用储电纳入同一接入限制；同一电量不得分别在供热与供冷台账重复计费或计碳。

\subsection{品质、可靠容量与地源账本}

舒适边界、设备和配电容量作为硬约束。可靠供热容量采用设备信用、可用率及储热持续供能条件，不能用储热额定kWh直接代替kW容量。储热在关键持续时间内能够提供的容量同时受功率和有效能量限制。建筑柔性只有在具有独立持续性依据的已启用模块下才能获得容量信用；未认证的RC调节量不自行抵扣设计峰值。

零失供及关键故障要求按服务合同和已实际测试情景记录。品质模块的启用取决于参数与实现，温湿度、结露、噪声和地源多年动态不能因存在接口即宣称验证通过。

地源热泵每个源分别核算地下取热$E_j^{\mathrm{ext}}$与排热$E_j^{\mathrm{rej}}$。供热取热扣除压缩机电耗，供冷排热包含制冷电耗；年度边界为
\begin{equation}
r_j^{\min}E_j^{\mathrm{rej}}\leq E_j^{\mathrm{ext}}
\leq r_j^{\max}E_j^{\mathrm{rej}}.
\label{eq:ground-ledger}
\end{equation}
该式按逐源和完整服务年检查。供暖季可能不启用年度比例约束，须标注其范围；年度比例通过也不证明多年地温稳定。

\clearpage
\section{投资调度求解与完整时序验证}

\subsection{成本、总碳与硬边界}

投资、网络选择与设备模式包含整数变量，运行状态和能流为连续变量，实际装配后形成MILP。成本目标写为
\begin{equation}
F_C=C_{\mathrm{cap}}^{\mathrm{ann}}+C_{\mathrm{energy}}
+C_{\mathrm{om}}+C_{\mathrm{carbon}}+C_{\mathrm{pen}},
\label{eq:cost-objective}
\end{equation}
其中资本年化参数、能源与运维价格、假设碳价和惩罚定义均来自数据字典。已启用的火用偏好成本归入惩罚项并单列。$F_E$为冷热服务与共享电力一致口径的系统总碳。可靠性、舒适和资源边界不与成本目标交换；同一费用或排放项只进入一个明确的组件台账。

\subsection{词典序端点与\texorpdfstring{$\varepsilon$}{epsilon}约束}

端点与碳帽扫描参考多目标约束法的组织原则\supcite{5}，具体容差和时限身份由本项目实际实现与运行记录确定。

模型提供两阶段词典序接口：对主目标$f_1$先求解，得到可行值$v_1$和上下界；再保留原约束，并增加
\begin{equation}
f_1\leq v_1+\delta_{\mathrm{abs}}+\delta_{\mathrm{rel}}|v_1|,
\label{eq:lexicographic}
\end{equation}
随后最小化次目标$f_2$。成本端点采用$(f_1,f_2)=(F_C,F_E)$，碳端点采用$(F_E,F_C)$；两个端点使用独立模型实例，记录两阶段状态、gap、容差和求解时间。第一阶段时间受限时，第二阶段仅在已获得可行值的容差边界内优化，不能称为已证明的全局词典序最优。

$\varepsilon$约束在统一可靠性与服务边界下求解
\begin{equation}
\min F_C,\qquad F_E\leq\varepsilon_E.
\label{eq:epsilon-constraint}
\end{equation}
两个扫描入口的目标分别记录：\path{experiments.py}的专用扫描将碳价设为零；\path{run_competition_frontier.py}保留情景配置的总目标，包括碳价成本及已启用的偏好项。因此赛事扫描结果不能标为纯粹的无碳价成本前沿。碳帽原值和数值容差后的有效值同时保存。检查每档状态、总碳及数值可行性后，对实际采样可行点剔除重复和支配解；有限采样集合不能写成完整精确Pareto前沿，也不声称实现了完整AUGMECON。

上述接口需由对应实验记录证明。单目标运行且未设置碳帽的manifest，不构成已完成$\varepsilon$扫描或词典序端点证明；供热原型的多目标运行与后述公开数学例题复现分别验收。

\subsection{代表期规划与投资冻结}

代表期保留气象、得热、舒适和服务状态的同源组合，记录实际日到代表日的映射、小时权重及覆盖审计。独立典型日难以表达跨期储热联系，代表期聚合须处理状态的时间联系\supcite{4}。本稿不声称实现该文的季节储能状态分解；状态边界单独给出，不能拼接不同年度输入以补齐所需小时，也不能将加权代表小时表述为完整连续运行。

先在代表期确定投资方案，保存容量、技术、网络、储热、接入和初末边界等记录。全年及供暖季运行沿用同一投资指纹和服务合同，新增或退役投资变量冻结；两种时序的范围和状态边界各自声明。完整时序由实际输入时间戳构成，运行层重新调度但不重新选择投资。

求解通过SciPy 1.15.3的\path{scipy.optimize.milp}接口调用HiGHS\supcite{6}。日志记录状态、原始与对偶界、相对gap、时限和数值审计；终止与可行性容差按接口及实际求解器选项解释\supcite{7}，不假定包装接口支持全部原生选项。时限内可行解保留其近似身份；通过运行审计仅说明所给投资在该时序内可运行，不说明全年度投资已联合优化到全局最优。

\subsection{独立数值与来源审核}

矩阵审计检查约束、变量上下界、整数可行性、目标重算及碳帽。建筑轨迹由独立RC复算核对采样温度、边界、冷热互斥和交付量。储热、逐源年度地源、资本一次计入和冷热总碳分别形成审计台账。

独立复算须使用可追溯的序列、配置和输入锚点，留存源文件与代码epoch指纹；只审查已保存矩阵的报告，不能宣称独立重建了全部设备、网络方程。来源审核另检查字段依据和变换身份。数值通过、来源通过和外部观测验证是不同证据，不能相互代替。

\clearpage
\section{验证证据、适用范围与对照}

\subsection{验证范围}

第三类场景以冷热共用设备、储热、建筑RC与配电边界建立原型服务模型。其规模、耗时、gap和审计残差属于软件验证；未经来源准入的容量、费用及排放属于情景数值，不形成真实工程改善结论。

主证据包括NASA产品全量来源复核、N2首轮与探索性比较，以及原作者多目标例题复现。投资运行的软件记录仍按代表规划、完整冷热服务年和供暖季分别报告。未来工程案例须关联通过来源审查的配置、规划与回放manifest及同条件基准。

\begin{table}[H]
\centering
\caption{最终案例所需的结果证据}
\begin{tabularx}{\textwidth}{@{}p{0.19\textwidth}p{0.28\textwidth}X@{}}
\toprule
记录 & 结果内容 & 核验条件\\
\midrule
规划 & 投资、代表期范围、状态与gap & 代码及来源版本确定，覆盖权重和建筑边界可追溯\\
年度 & 完整冷热服务、费用、总碳与地源 & 连续逐时、投资冻结、资本一次计入、逐源年度账本\\
供暖季 & 该期间交付热量、状态与运行费 & 独立时钟和边界；不与年度资本或全年地源条件混用\\
建筑审计 & 温度、RC残差、冷热互斥及交付量 & 轨迹同源、时段末采样范围明确；观测精度另验证\\
公平基准 & 同服务条件下的差值 & 同天气、建筑、舒适、端点和成本碳口径，来源准入通过\\
\bottomrule
\end{tabularx}
\end{table}

\subsection{公平比较与外部证据}

基准和推荐方案采用相同气象、建筑参数、服务需求、舒适与初末状态、失供要求及会计边界。技术限制或容量调整必须预先声明，不能通过降低服务质量、删除冷需求或移除资本使某方案显得更优。基准未完成同条件核验时，仅展示程序状态，不给出改善结论。

外部RC观测的留出验证与投资调度比较分开报告。只有真实观测、参数辨识划分和留出预测全部可追溯时才计算外部误差；无法获得观测的参数继续保留假设身份。模型求解记录及程序审计不作为企业认可、鉴定或现场测量的替代材料。

\Needspace{8\baselineskip}
\subsection{独立外部RC适用性验证}

采用Fraunhofer/IEA EBC Annex 71德国N2主实验发布数据\supcite{8}，对单一1R1C结构作外部适用性诊断。数据版本2、DOI与CC BY-SA 4.0许可可查，全部11988条连续记录按时间分为训练、验证和测试段，数量分别为7192、2397和2399。独立复算核对来源文件及162项指标，首轮结果已获有限方法引用准入。

该资料是发布方处理后的测量数据。短缺失温度插值及部分功率、流量置零没有逐行质量旗标，无法剔除全部受处理样本，也不能排除源处理跨划分信息。选择N2排除其他房屋明确的模拟补热，不等于获得未经处理的原始传感器数据。

十个房间110~cm测点温度按地板面积加权形成派生观测，不称全屋体积均温。加热及内部热源功率由W转换为kW，电热按1倍换算为显热是声明假设。验证采用区间平均观测方程；室内与功率十分钟聚合窗口、标签相位及计量不确定度尚未完整确认，方程代数精确不证明秒级过程精确。太阳得热采用声明代理，窗面积与SHGC乘积不作为认证全屋严格热增益上界。

训练、验证和测试按连续时间划分，以独立数据检验预测表现\supcite{3}。本次参数与起点选择仅依据训练段；内部方案文件不单独证明不可变事前注册。一步预测使用上一实测区间温度，连续开环仅初始化一次，分段开环仅在各划分段起点初始化；持续温度基准采用同样的信息和初始化规则。

\begin{table}[H]
\centering
\caption{N2首轮与24~h损失探索的测试段RMSE（$^\circ$C）}
\begin{tabularx}{\textwidth}{@{}Xrrr@{}}
\toprule
预测模式与初始化 & 首轮RC & 后续探索RC & 持续性基准\\
\midrule
一步：上一实测区间温度 & 0.085979 & 0.091334 & 0.091964\\
连续开环：全序列初始化一次 & 1.885259 & 7.425830 & 0.790172\\
分段开环：各划分段初始化一次 & 1.155136 & 6.170260 & 0.626294\\
\bottomrule
\end{tabularx}
\end{table}

首轮一步误差略低于基准，两种开环均高于基准，揭示长时漂移。后续采用24~h日块训练损失的探索沿用同一划分，测试开环进一步恶化，太阳代理参数达到上界；训练改善未转化为留出改善，不能称新独立验证。完整重拟合独立复现首轮与后续参数及指标。结果仅支持已知天气和热输入下有限的条件单步适用性；德国拟合H/C不迁移到武汉，也不证明中国工程效果。

\clearpage
\subsection{原作者多目标算例复现}

采用Mavrotas作者公开报告的双目标连续LP例题\supcite{5}：
\begin{equation}
\max f_1=x_1,\qquad \max f_2=3x_1+4x_2,
\end{equation}
\begin{equation}
0\leq x_1\leq20,\quad0\leq x_2\leq40,\quad5x_1+4x_2\leq200.
\end{equation}
非负域依据原图2解释，其余系数按原式录入。复现在项目LinearMILP数值层上另用两阶段包装：先最大化主目标，再等式固定实际所得主目标值并优化次目标。两端点$(f_1,f_2)$为$(20,160)$与$(8,184)$，与作者表2一致。

由实际端点的$f_2$范围四等分，逐档施加$f_2\geq\varepsilon$并分层求解。表\ref{tab:published-epsilon}给出实际结果，五点唯一且互不支配，与独立有理数多边形解析界一致。
\begin{table}[H]
\centering
\caption{公开数学例题的实际分层$\varepsilon$约束解}
\label{tab:published-epsilon}
\begin{tabular}{rrrrr}
\toprule
$\varepsilon$ & $x_1$ & $x_2$ & $f_1$ & $f_2$\\
\midrule
160 & 20 & 25 & 20 & 160\\
166 & 17 & 28.75 & 17 & 166\\
172 & 14 & 32.5 & 14 & 172\\
178 & 11 & 36.25 & 11 & 178\\
184 & 8 & 40 & 8 & 184\\
\bottomrule
\end{tabular}
\end{table}

\begin{figure}[H]
\centering
\begin{tikzpicture}[x=0.5cm,y=0.13cm,font=\small]
\draw[arrow] (-1,-2) -- (14,-2) node[right] {$f_1$（最大化）};
\draw[arrow] (-1,-2) -- (-1,28) node[above] {$f_2$（最大化）};
\foreach \x/\lab in {0/8,3/11,6/14,9/17,12/20}{
  \draw (\x,-2) -- (\x,-2.7) node[below] {\lab};
}
\foreach \y/\lab in {0/160,6/166,12/172,18/178,24/184}{
  \draw (-1,\y) -- (-1.2,\y) node[left] {\lab};
}
\draw[thick,draw=paperGreenDark] (0,24) -- (12,0);
\foreach \x/\y in {0/24,3/18,6/12,9/6,12/0}{
  \fill[paperGreenDark] (\x,\y) circle[radius=2pt];
}
\node[align=left] at (8,23) {实点：实际LP解\\实线：独立解析前沿};
\end{tikzpicture}
\caption{公开数学例题的五个采样点与解析前沿}
\end{figure}

含资源上限180/220的声明边界检验，共实际调用23次LP：22次最优，1次预期不可行；最大矩阵行、变量界违反及目标重算误差均为0。累计求解0.017743~s属于此小例题，不代表工程求解速度；未暴露的LP gap保留null。复现未调用\path{HeatPlanningModel.solve_lexicographic}，不属于完整AUGMECON或供热工程收益验证。项目最终全量回归628项通过，耗时67.55~s，只作为软件检查记录。


\clearpage
\section{经济碳评价、边界与局限}

\subsection{费用与排放解释}

全年评价将规划资产的年化资本与全年运行费用核对，资本只计一次；冻结投资后，若资本已在年度模型目标中计入，后处理不得再加一遍。供暖季仅反映该期间运行，不能按小时比例外推全年费用，也不能与含全年冷热资本的目标直接比较。

费用台账区分资本年化、能源、运维、假设碳价及惩罚项。惩罚项不是企业现金支付；假设碳价成本不据此表述为企业已实际缴纳的碳费。总目标、含假设碳价的情景费用和不含该项的费用分别展示，分项之和与求解目标交叉核对。

系统总碳覆盖供热、供冷及共享电力的统一口径。若程序分别输出供热台账和供冷台账，应按其定义合成总量，不重复计算共享用电，也不能只取供热部分作为全年系统排放。单位冷热服务指标须使用对应交付热量和冷量作分母；只有两方案服务、输入和会计口径一致时才计算差值和改善比例。

\subsection{能力与适用范围}

精确1R1C是给定小时内输入保持条件下的状态方程离散，不等于建筑已经实测校准。独立RC审计认证时段末采样状态、边界和账本一致性，不据此声称房间级温度、实测误差或所有时段内连续舒适已经验证。

代表期边界会影响冷热需求与容量。每天强制恢复同一温度可能产生额外的恢复冷热量，应结合完整时序及边界敏感性检查后评价规划结果。代表期投资方案的全年回放检验其运行可达性，不证明该方案是全年联合投资模型的全局最优。

年度地源取排热账本只检验本年度能量比例，不证明多年地温稳定。聚合网络不替代真实GIS、连续水力、配电潮流或施工图。热湿、结露、声学、地质和多年动态需独立资料与模块支持；未启用或未通过专业复核的内容不列为成果。C\&CG、Benders分解、自动工程割和在线MPC亦不作为本稿已验证的求解方法。

\clearpage
\section{结论与复现}

本稿建立三场景接口下的源网荷储、建筑RC、品质容量和成本碳台账，采用代表期规划、冻结投资完整时序运行和独立审计衔接方案生成与验证。方法将输入来源审核与数值验收分开：真实求解不证明输入真实，数值自洽不等于实测校准。

气象来源复核认证四个产品字段，N2比较揭示条件单步与开环表现的差异，后续探索未改善留出开环表现。公开数学例题的两个端点、五个采样点和声明边界检验与解析解一致，支持数值层方法复现。三类证据分别提供来源追溯、模型适用性诊断和软件方法检查，工程经济改善仍需真实参数、报价及同条件基准。企业认可与鉴定依据另行提供。

复现资料包含配置、逐时输入、预处理说明、求解日志、manifest、投资记录、运行轨迹及独立审计。每个结果表均应关联其实际运行目录和哈希；论文只展示必要公式与伪代码，完整程序保留在工程中。

\clearpage
\begin{center}{\heiti\sihao 参考资料}\end{center}
\addcontentsline{toc}{section}{参考资料}
\begin{refsectionflush}
\refentry{1}{中国建筑科学研究院有限公司. 面向供热格局变革的城镇供热源荷匹配与规划寻优模型构建比赛方案{[}R{]}. 2026，第4页.}
\refentry{2}{Bacher P, Madsen H. Identifying suitable models for the heat dynamics of buildings{[}J{]}. Energy and Buildings, 2011, 43(7): 1511--1522. DOI: \url{https://doi.org/10.1016/j.enbuild.2011.02.005}.}
\refentry{3}{Madsen H, Bacher P, Bauwens G, et al. Thermal performance characterisation using time series data -- statistical guidelines{[}R{]}. IEA EBC Annex 58, Subtask 3, Part 2. KU Leuven, 2016. \url{https://iea-ebc.org/Data/publications/EBC_Annex_58_Final_Report_ST3b.pdf}.}
\refentry{4}{Kotzur L, Markewitz P, Robinius M, Stolten D. Time series aggregation for energy system design: Modeling seasonal storage{[}EB/OL{]}. arXiv:1710.07593v2, 2018. \url{https://arxiv.org/abs/1710.07593}. DOI: 10.1016/j.apenergy.2018.01.023.}
\refentry{5}{Mavrotas G. Effective implementation of the $\varepsilon$-constraint method in Multi-Objective Mathematical Programming problems{[}J{]}. Applied Mathematics and Computation, 2009, 213(2): 455--465. DOI: 10.1016/j.amc.2009.03.037. 作者配套方法报告：\url{https://www.gams.com/modlib/adddocs/epscm.pdf}.}
\refentry{6}{SciPy developers. scipy.optimize.milp -- SciPy 1.15.3 Reference Manual{[}EB/OL{]}. [2026-10-02]. \url{https://docs.scipy.org/doc/scipy-1.15.3/reference/generated/scipy.optimize.milp.html}.}
\refentry{7}{HiGHS developers. List of options -- HiGHS Documentation{[}EB/OL{]}. 开发版文档，[2026-10-02]. \url{https://ergo-code.github.io/HiGHS/dev/options/definitions/}.}
\refentry{8}{Kersken M, Strachan P. Twin House Experiment IEA EBC Annex 71 Validation of Building Energy Simulation Tools -- Specifications and dataset{[}DS{]}. Fraunhofer Fordatis, version 2, 2020-06-23. CC BY-SA 4.0. \url{https://doi.org/10.24406/fordatis/76.2}.}
\refentry{9}{NASA POWER. Hourly API{[}EB/OL{]}. [2026-10-02]. \url{https://power.larc.nasa.gov/docs/services/api/temporal/hourly/}. 本次响应版本v2.10.2，原始响应、查询及逐字段比对记录随复现资料保存.}
\end{refsectionflush}

\clearpage
\section*{附录A\quad 场景接口与审计式}
\addcontentsline{toc}{section}{\texorpdfstring{附录A\quad 场景接口与审计式}{附录A 场景接口与审计式}}

\subsection*{场景接口}
第一类加载既有资产、热源替代、退役及温位迁移；第二类加载选址、用户分配、燃料和微网；第三类加载冷热服务、建筑RC、设备共用容量与配电。只有具有字段依据、实现和实际求解记录的模块才列为验证成果。当前完整案例证据针对第三类场景组织，第一、二类不沿用历史数字冒充本轮验证。

\subsection*{关键审计}
逐时热量平衡残差记为$r_{n,g,t}^{Q}$，储热状态残差为$r_{s,t}^{E}$，建筑状态残差为$r_{b,t}^{\theta}$：
\begin{align}
\max_{n,g,t}|r_{n,g,t}^{Q}|&\leq\varepsilon_Q,\\
\max_{s,t}|r_{s,t}^{E}|&\leq\varepsilon_S,\\
\max_{b,t}|r_{b,t}^{\theta}|&\leq\varepsilon_T.
\end{align}
容差按物理单位及保存的数值审计声明，状态残差不与温度观测误差混同。另检查采样舒适、初末状态、冷热互斥、设备和配电上界、整数与目标重算；地源比例依据式\eqref{eq:ground-ledger}仅在启用的年度范围审计。

来源记录、服务合同、代码版本、规划投资和回放输入指纹一致后，仍需区分来源准入与数值验收。只通过数值验收的假设数据案例保持程序验证身份。未执行的压力测试、$\varepsilon$扫描、观测留出及专业工程复核，不作为通过项目列入结果。


\clearpage
\section*{附录B\quad 求解与验证伪代码}\label{ux9644ux5f55-a-ux4e09ux573aux666fux63a8ux8350ux914dux7f6eux6a21ux677f}
\addcontentsline{toc}{section}{\texorpdfstring{附录B\quad 求解与验证伪代码}{附录B 求解与验证伪代码}}

\noindent\textbf{输入：}字段来源记录、配置、逐时数据、服务合同、投资候选、边界条件、求解时限与容差；启用多目标时另给定碳帽与词典序容差。\textbf{输出：}规划与运行记录、成本总碳台账、审计结果及证据身份。方法伪代码与已经执行的实验记录分别核验。

\begingroup
\small
\setlength{\LTpre}{4pt}
\setlength{\LTpost}{4pt}
\renewcommand{\arraystretch}{1.15}
\begin{longtable}{@{}p{0.07\textwidth}@{\hspace{0.03\textwidth}}p{0.90\textwidth}@{}}
\toprule
步骤 & 规划、运行与证据准入\\
\midrule
\endfirsthead
\toprule
步骤 & 规划、运行与证据准入（续）\\
\midrule
\endhead
\bottomrule
\endfoot
01 & 审核每个输入字段的来源、单位和时间轴；保存原始文件、变换规则及指纹。\\
02 & \textbf{如果}工程字段来源未通过，结果身份限于程序验证/情景示例；不得输出真实工程绩效结论。\textbf{如果}时间轴或单位不合法，停止相应案例。\\
03 & 生成候选和代表期，检查覆盖权重、建筑初末状态及储热状态；装配实际启用的源网荷储、RC与容量约束。\\
04 & \textbf{如启用多目标}，用独立模型实例求两个词典序端点：先求主目标，再用显式容差约束主目标并求次目标；保存两阶段状态与上下界。\\
05 & \textbf{对每个}启用的碳帽求成本最小解，核对系统总碳及数值可行性；单目标案例使用其已声明目标，不视为完成$\varepsilon$扫描。\\
06 & \textbf{如果}无可行投资方案，保存状态并结束该案例；\textbf{否则}，保存全部投资和服务合同指纹并冻结投资。\\
07 & \textbf{对完整服务年和独立供暖季}，在连续逐时输入上重求运行，按各自范围核对状态、总碳、资本与时钟。\\
08 & \textbf{对每个建筑和时段}，独立复算RC状态；检查采样舒适、初末边界、冷热互斥和交付量，另核对储热及矩阵审计。\\
09 & 年度逐源检查地源取排热账本；季节记录其约束启用范围。核对资本只计一次、费用分项及总碳口径。\\
10 & \textbf{如有公平基准}，先核对同天气、建筑、舒适和服务条件，再计算差值；条件不一致时停止优越性比较。\\
11 & \textbf{如果}来源、求解或审计未达到相应准入条件，保留问题与阶段身份；\textbf{否则}，输出该范围内可追溯的结果。\\
12 & 遍历已配置案例后停止；时间受限解保留gap和上下界，不自动生成工程割，不把未执行的步骤写为实验成果。\\
\end{longtable}
\endgroup
\end{document}
